\section{Technique Application According to Course Guideline}

\subsection{Guideline Alignment}
The course guideline at
\url{https://datalab.vn/articles/2021-09/IT5425-quan-tri-du-lieu-va-truc-quan-hoa}
asks each project to explain which visualization techniques are used and where
they appear in the final artifact. We use the course slides and the distilled
note file \code{Slide data viz/2024/notes/00\_distilled\_knowledge.md} as the
reference, then connect the concepts directly to the three final dashboard
pages.

\begin{table}[h]
\centering
\caption{Course knowledge used in the final dashboard}
\begin{tabularx}{\textwidth}{p{0.24\textwidth}p{0.34\textwidth}X}
\toprule
\textbf{Course knowledge} & \textbf{How we used it} & \textbf{Dashboard evidence} \\
\midrule
Visual encoding &
We map categories to bar positions and quantitative values to bar length. Color
is reserved for churn/risk emphasis. &
Overview KPI page, selected-feature churn-rate chart, and Random Forest
importance bars. \\
Graphical perception &
We rely on common baselines, sorted categories, and preattentive orange
highlights instead of dense legends. &
Churn driver page and dynamic risk chart rank higher-risk categories first. \\
Table-data visualization &
We aggregate customer-level rows into feature groups, rates, counts, and model
summary tables. &
Driver detail table, confusion matrix, split table, and risk-by-feature table. \\
Figure design &
We avoid decorative 3D effects, keep direct labels, state denominators, and use
a gray--orange--blue palette for accessibility. &
All final screenshots use flat charts, KPI cards, labels, and consistent risk
colors. \\
Interaction &
We use slicers and view changes so one dashboard page can answer several
feature-level questions. &
Feature filters on the driver page and the single-select dynamic risk chart on
the prediction page. \\
Storytelling &
We order the dashboard from context to drivers and then to model-assisted
retention priority. &
Page 1 gives the churn baseline, page 2 explains drivers, and page 3 supports
retention targeting. \\
\bottomrule
\end{tabularx}
\end{table}

\subsection{Chapter 1: Overview of Data Visualization}
\paragraph{Techniques / principles applied.}
We use visualization for both exploration and explanation. The EDA views help
the group compare churned and non-churned customers, while the final dashboard
explains which customer groups are above the baseline churn rate.

\paragraph{How applied in the dashboard.}
The overview page starts with total customers, churn count, churn rate, and risk
distribution. This gives the audience a baseline before the report moves into
feature-level churn drivers and prediction outputs.

\paragraph{Notes / adjustments.}
The raw dataset is too large for direct row-level inspection during
presentation, so the dashboard uses aggregated segment views and summary cards.

\subsection{Chapter 2: Visual Models and Visual Encoding}
\paragraph{Techniques / principles applied.}
The selected variables have different measurement levels: contract type is
nominal, satisfaction and payment groups are ordinal/discrete, and monthly
charges or usage are quantitative. We therefore use categorical axes for groups
and bar length for counts, churn rates, feature importance, and predicted risk.

\paragraph{How applied in the dashboard.}
The selected-feature churn-rate chart maps feature categories to bars and
within-category churn rate to bar length. The dynamic risk chart uses the same
logic: the slicer selects a business feature, the axis shows its categories, and
the bar length shows average predicted churn probability.

\paragraph{Notes / adjustments.}
We avoid stacking unrelated risk metrics into one visual. Churn rate, customer
share, and average predicted probability are kept as separate metrics because
they answer different questions.

\subsection{Chapter 3: Graphical Perception}
\paragraph{Techniques / principles applied.}
We use encodings that people compare accurately: position, length, ordering, and
clear baselines. Color is used to guide attention, not to carry the whole
meaning.

\paragraph{How applied in the dashboard.}
The driver page places feature-level charts and detail tables together, so the
viewer can compare category risk and segment size in the same context. The
dynamic risk chart sorts categories by average predicted churn probability, so
the highest-risk category appears first.

\paragraph{Notes / adjustments.}
The gray--orange--blue palette helps draw attention to risk while staying more
readable than red-green contrast. Labels, sorting, and chart titles remain
available for viewers who cannot rely on color.

\subsection{Chapter 4: Visualization for Table and Multidimensional Data}
\paragraph{Techniques / principles applied.}
The project is mainly a table-data visualization problem. We transform
customer-level records into grouped features, rates, counts, feature-importance
tables, confusion matrices, and risk summaries.

\paragraph{How applied in the dashboard.}
The churn-driver page compares grouped categories against the overall churn
baseline. The prediction page uses the long-format \code{RiskByFeature} table so
the same bar chart can switch between contract type, satisfaction, complaints,
service calls, late payments, and other selected features.

\paragraph{Notes / adjustments.}
Churn-rate bars and average-probability bars do not sum to 100\%. Each bar is a
within-category metric. Only composition charts, such as share of churned
customers, are expected to sum to a whole.

\subsection{Chapter 5 and Chapter 7: Graph and Map Visualization}
\paragraph{Techniques / principles applied.}
Graph and map techniques are not central to this dataset because the available
fields do not contain network edges or geographic locations.

\paragraph{How applied in the dashboard.}
The final dashboard therefore focuses on table-data visualization rather than
adding graph layouts or maps that the data cannot support.

\paragraph{Notes / adjustments.}
If later data adds regions, stores, or customer relationships, this section can
be extended with map encoding or graph-layout choices.

\subsection{Chapter 6: Principles of Figure Design}
\paragraph{Techniques / principles applied.}
We apply proportional ink, clear titles, readable labels, consistent colors, and
simple layouts. The final pages avoid 3D effects and unnecessary decoration
because the main task is comparison.

\paragraph{How applied in the dashboard.}
The overview page uses KPI cards for scale, the driver page uses bars and
tables for segment comparison, and the prediction page combines model metrics
with the dynamic risk chart. The gray--orange--blue palette gives a stable
visual language: orange for warning or higher risk, blue for the main selected
comparison, and gray for context.

\paragraph{Notes / adjustments.}
The dashboard still has room for typographic polish, but the report focuses on
the design decisions that materially affect interpretation: denominator clarity,
aggregation choice, sorting, and color accessibility.

\subsection{Chapter 8: Interactive Visualization}
\paragraph{Techniques / principles applied.}
The dashboard uses filtering, selection, and view changes. Interaction is used
to support comparison rather than to hide the baseline or replace explanation.

\paragraph{How applied in the dashboard.}
The most advanced interaction is the Power BI dynamic risk chart. It uses a
single-select slicer based on \code{RiskByFeature[feature\_display]}; the bar
chart then shows \code{RiskByFeature[category]} and the Average of
\code{RiskByFeature[avg\_churn\_probability\_pct]}. Selecting contract type
changes the bars to month-to-month, one-year, and two-year contracts; selecting
late payments changes the bars to late-payment groups.

\paragraph{Notes / adjustments.}
This chart is not a risk distribution. It compares average predicted churn
probabilities, so the bars are not expected to add up to 100\%. Using Average
instead of Sum prevents probability values from being interpreted incorrectly.

\subsection{Chapter 9: Storytelling with Data}
\paragraph{Techniques / principles applied.}
The narrative moves from context to diagnosis and then to action. This matches
the retention use case: managers first need the churn baseline, then the main
drivers, and finally a way to prioritize customers.

\paragraph{How applied in the dashboard.}
The three final pages follow that order. Page 1 frames the overall churn
problem, page 2 explains feature-level churn drivers, and page 3 connects model
outputs with retention priority through predicted churn risk and the confusion
matrix.

\paragraph{Notes / adjustments.}
The final wording avoids treating the model as a perfect classifier. The model
supports prioritization, while the dashboard and report provide the business
context needed to interpret the predictions.
